% !TEX root = main.tex

\section{Introduction}
\label{sec:intro}

Between 2016 and 2025, large-capitalisation United States technology equities
traded through several distinct macro-financial regimes: an extended period of
near-zero policy rates, the COVID-19 market dislocation of March 2020, the most
rapid monetary tightening cycle since the early 1980s beginning in 2022, and the
partial normalisation of rates that followed. These episodes produced pronounced
volatility clustering, heavy-tailed daily returns, and repeated shifts in the
joint distribution of returns and macroeconomic state variables, features long
documented as stylised facts of asset returns \citep{cont2001empirical}. Such
conditions are hostile to forecasting models estimated under an assumption of
covariance stationarity, and they reproduce at the level of a portfolio a
persistent finding of the empirical finance literature: predictive relations
that appear stable in sample frequently deteriorate out of sample
\citep{welch2008comprehensive, timmermann2008elusive}. As \citet{west2020bayesian}
argues, structure uncertainty then becomes as consequential as parameter
uncertainty, since the analyst must decide not only what the coefficients are
but which model, if any, remains adequate as the regime changes.

There are sound theoretical grounds for conditioning return forecasts on
macro-financial information. Multifactor asset pricing theory implies that
innovations in systematic state variables are priced in equity returns
\citep{chen1986economic}. Empirically, unanticipated changes in the federal funds
rate induce large and statistically significant equity market responses
\citep{bernanke2005explains}; option-implied volatility, summarised by the CBOE
Volatility Index, is contemporaneously and negatively associated with returns
\citep{whaley2000investor}; and comovement with the broad market index remains
the dominant source of systematic variation in individual stock returns. Recent
data-driven studies confirm that volatility and uncertainty measures rank among
the most influential predictors of index returns \citep{latif2025integrating},
and that volatility transmission across markets is itself time varying
\citep{agrawal2026volatility}. The unresolved methodological question is
therefore not whether such variables carry information, but how they should enter
the forecasting model and how the resulting predictive uncertainty should be
quantified.

This paper approaches that question constructively. Most practical uses of a
return forecast, such as setting a risk limit, sizing a position or pricing
an option, require a predictive \emph{distribution} rather than a point estimate,
yet the forecasting literature is dominated by evaluations of conditional-mean
accuracy. We develop a model that supplies a calibrated distribution, and we
establish which of its components does the work.

The specification is a hierarchical Bayesian regression with four features. It
treats a GARCH(1,1)-$t$ conditional standard deviation as a \emph{known}
observation-level scale, so the model does not have to relearn volatility
clustering and a single scalar measures whether that scale is correctly sized.
It pools intercepts across assets, which allows an asset absent from estimation
to be forecast by integrating its intercept over the population distribution. It
places a horseshoe prior \citep{carvalho2010horseshoe} on the mean coefficients,
so that predictor selection occurs inside the posterior with uncertainty
attached rather than as a screening step beforehand. And it uses Student-$t$
innovations with a scale calibrated on validation data, which separates
tail behaviour from scale. Estimation is by blocked Gibbs sampling with every
conditional in closed form.

The evaluation is a controlled comparison on a balanced panel of eight large-cap
technology stocks over 2016--2025, with the effective federal funds rate, the
ten-year Treasury yield, the VIX and a market proxy as macro-financial
conditioning information. Competing specifications (naive baselines, pooled
ridge, a pooled random forest, three GARCH-family volatility models and four
sequence architectures) are estimated on an identical feature set and scored on
a strictly held-out window spanning the tightening cycle and its aftermath.
Point accuracy is assessed by root mean squared error and out-of-sample $R^2$;
probabilistic adequacy by logarithmic score, continuous ranked probability score
and interval calibration. Every comparison is tested by Diebold--Mariano rather
than ranked.

We report four findings. First, the proposed model attains the best logarithmic
score and CRPS among all specifications considered, improving on a
ridge-plus-GARCH benchmark by $0.126$ nats ($p<0.00001$), and reduces mean
absolute interval-coverage error to $0.022$ from $0.030$. Second, partial pooling
yields calibrated intervals for an asset with no estimation history at all,
attaining $87\%$ coverage at the $90\%$ nominal level and a logarithmic score
within $0.005$ nats of the same model estimated with that asset included. Third,
decomposing the gain in predictive density over a null forecast locates its
source: the conditional variance contributes $54\%$ and the shape of the
predictive distribution a further $45\%$, while the conditional mean contributes
$0.7\%$, detectable at $p=0.029$ and economically negligible. Fourth,
consistent with \citet{welch2008comprehensive} and \citet{timmermann2008elusive},
no specification improves on a per-ticker historical average for the conditional
mean, a result that survives disaggregation to individual assets and correction
for multiple testing, and which is precisely why modelling effort at this horizon
is better directed at the second moment and the distributional shape.

We also document a failure mode that is silent rather than merely severe.
Standardising non-stationary macroeconomic levels on a fixed estimation origin
sends held-out inputs far outside the range seen in estimation, and every model
with an unbounded output head extrapolates along them, producing out-of-sample
$R^2$ as low as $-10\,425\%$ with no exception raised and no implausible training
diagnostic. Because the defect is detectable only by auditing predictor
distributions across splits and the scale of the predictions themselves, both
audits are released with the code.

The remainder of the paper reviews the relevant literature, describes the data
and methodology, reports the empirical results, and concludes with
recommendations for practice.
